\documentclass[11pt]{article}

\usepackage{amsmath}
\usepackage{amssymb}
\usepackage{array}
\usepackage{booktabs}
\usepackage{bm}
\usepackage{geometry}
\usepackage{hyperref}
\usepackage{xcolor}

\geometry{a4paper, margin=1in}
\hypersetup{colorlinks=true, linkcolor=blue!50!black, citecolor=blue!50!black, urlcolor=blue!50!black}

\newcommand{\stream}{STREAM}
\newcommand{\uce}{UCE}
\newcommand{\cre}{cCRE}
\newcommand{\E}{\mathbb{E}}
\newcommand{\R}{\mathbb{R}}
\newcommand{\N}{\mathcal{N}}
\newcommand{\dd}{\mathrm{d}}

\title{Proposal: a learned stochastic developmental clock for STREAM}
\author{}
\date{}

\begin{document}

\maketitle

\begin{abstract}
Developmental stage is an imperfect proxy for the biological progress of an
individual cell. A population collected at one stage can contain cells moving
at different rates, yet precomputed pseudotime would introduce a heuristic
coordinate outside the dynamics model. We propose instead to augment
\stream{} with a learned stochastic clock. A UCE-conditioned scalar process
controls how far a cell travels along a normalized, gene-valued regulatory
velocity direction, while a separate residual diffusion accounts for
variation not explained by progress rate. The clock is learned from unpaired
endpoint populations by differentiable rollout and requires neither
pseudotime labels nor absolute-time conditioning. This document distinguishes
the clock from the current gene-wise score SDE, defines identifiable
parameterizations and positive increments, describes fitting and validation,
and specifies ablations that test whether variable developmental progress is
responsible for the current persistence-level forecasts.
\end{abstract}

\section{Motivation}

The causal \stream{} model defines an autonomous expression-space field
\begin{equation}
  \frac{\dd x}{\dd t}
  = v_\theta\!\left(E(x),\{Z_g\}_{g=1}^{G}\right),
  \label{eq:current-ode}
\end{equation}
where $x\in\R^G_{\geq 0}$ is gene expression, $E$ is the frozen \uce{}
encoder, and $Z_g$ contains AlphaGenome-derived regulatory tokens linked to
gene $g$. The output remains one velocity per gene. Absolute developmental
time is not an input; \uce{} is recomputed from the current expression state
during rollout.

Equation~\ref{eq:current-ode} nevertheless assumes that the experimental time
increment corresponds to the same amount of biological progress for every
cell. That assumption is unlikely to hold exactly. Cells sampled at one stage
may be delayed, rapidly differentiating, cycling, quiescent, or already
committed to different fates. Diffusion pseudotime, Palantir, and related
methods infer useful orderings from expression geometry
\cite{haghverdi2016dpt,setty2019palantir}, while dynamical RNA-velocity models
can infer gene-specific latent times from splicing kinetics
\cite{bergen2020scvelo}. However, supplying one of these coordinates to
\stream{} would make the model depend on a separately chosen graph, root,
feature space, and smoothing procedure.

The desired alternative is to make progress itself latent and learn it jointly
with the expression dynamics. Experimental time then specifies the duration
over which a cell may progress, rather than directly fixing the distance it
must travel through expression space.

\section{Why the current SDE is not already a clock}

The coupled score-flow model currently supports an SDE of the form
\begin{equation}
  \dd X_t
  = \left[v_\theta(X_t)+\kappa\Sigma s_\theta(X_t)\right]\dd t
  + \sqrt{2\kappa}\,\Sigma\dd W_t,
  \label{eq:current-sde}
\end{equation}
where $\Sigma=\operatorname{diag}(q)$ contains fixed gene scales and
$W_t\in\R^G$ has one Brownian coordinate per gene. The corresponding regular
stochastic-interpolant target uses
$\epsilon\sim\N(0,I_G)$, also a $G$-vector. This noise can cause some particles
to move farther than others, so it may incidentally create variation in
apparent progress. It does not, however, encode progress explicitly:
\begin{itemize}
  \item the Brownian increment has one component per gene and need not align
        with the developmental velocity;
  \item much of its variation can be perpendicular to the learned trajectory;
  \item the score correction is designed to control population marginals, not
        to estimate a cell-specific developmental rate; and
  \item a wider endpoint distribution can improve Sinkhorn distance without
        producing a coherent or interpretable ordering.
\end{itemize}

A developmental clock should instead perturb all genes coherently along the
same local velocity direction. Residual gene-wise diffusion may still be
useful, but it should be represented and evaluated separately.

The coordinate transform introduced below also simplifies the non-clock model.
With $Y=Q^{-1}X$, a deterministic count-space field and SDE transform as
\begin{align}
  \dd X&=b_X(X)\dd t+QA_Y(Y)\dd W,
  &
  \dd Y&=Q^{-1}b_X(QY)\dd t+A_Y(Y)\dd W.
  \label{eq:general-coordinate-sde}
\end{align}
Thus the existing $Q\dd W$ perturbation becomes identity-scaled in $Y$.
Likewise, ordinary flow matching in $Y$ is exactly count-space flow matching
with each velocity divided by its gene scale. This coordinate choice applies
to deterministic CFM, regular score-flow, explicit SDE rollout, and the latent
clock; it is not specific to pseudotime.

\section{Model}

\subsection{Notation and dimensions}

Let $G$ be the number of modeled genes, $D$ the \uce{} state dimension, $K$
the maximum number of regulatory tokens per gene, $C$ the CRE-token embedding
dimension, and $r$ the optional residual diffusion rank. The main variables
have the following dimensions:
\begin{table}[ht]
\centering
\small
\begin{tabular}{@{}p{0.22\textwidth}p{0.24\textwidth}p{0.46\textwidth}@{}}
\toprule
Variable & Dimension & Meaning \\
\midrule
$X_t,q$ & $\R^G$ & Count-like expression state and positive gene scales \\
$Y_t,u,d,\widehat d$ & $\R^G$ & Gene-scaled state, direction, or displacement \\
$E(X)=E(QY)$ & $\R^D$ & Frozen \uce{} state after returning to expression coordinates; currently $D=1280$ \\
$Z$ & $\R^{G\times K\times C}$ & Padded gene-by-regulatory-token representations \\
$S_t,a_\phi,\sigma_\phi$ & Scalars & Progress, expected clock rate, and clock variability \\
$W_t^S,z,\widehat z$ & Scalars & Clock Brownian motion, sampled clock noise, and predicted clock noise \\
$W_t^\perp$ & $\R^r$ & Optional residual Brownian motion \\
$G_\psi(x)$ & $\R^{G\times r}$ & Residual expression-diffusion loading matrix \\
$\epsilon,W_t$ & $\R^G$ & Gene-wise noise in the regular score-flow model \\
$\tau,\gamma(\tau),\rho$ & Scalars & Interpolation coordinate, bridge schedule, and fixed bridge amplitude \\
\bottomrule
\end{tabular}
\caption{Shapes of the variables used in the stochastic-clock model.}
\label{tab:dimensions}
\end{table}

The vector $q=(q_1,\ldots,q_G)\in\R^G_{>0}$ contains fixed training-only
per-gene expression scales estimated from endpoint counts with a nonzero
floor. Define $Q=\operatorname{diag}(q)$ and use the model coordinate
\begin{equation}
  Y=Q^{-1}X,\qquad X=QY.
  \label{eq:coordinate-transform}
\end{equation}
All dynamics, losses, direction normalization, and nonnegative projection are
performed in $Y$. The fixed $q$ appears only at system boundaries: $X=QY$ is
passed to \uce{}, PCA/Sinkhorn evaluation, and biological reporting.
Gene-scaled coordinates remain nonnegative and avoid the library-size state
needed to invert log normalization. The latent progress $S$ is dimensionless,
$a_\phi$ has units of inverse experimental time, and $\sigma_\phi$ has units
of inverse square-root time. The bridge coordinate $\tau$, bridge amplitude
$\rho$, and clock draw $z$ are dimensionless. Crucially, $z\in\R$ is one
scalar shared across all genes for one cell and path, unlike the regular
score-flow noise $\epsilon\in\R^G$.

\subsection{Latent stochastic clock}

Introduce a scalar biological-progress process $S_t$ and write
\begin{align}
  \dd S_t &= a_\phi(Y_t)\,\dd t + \sigma_\phi(Y_t)\,\dd W_t^{S},
  \label{eq:clock-sde}\\
  \dd Y_t &= u_\theta(Y_t)\,\dd S_t
             + G_\psi(Y_t)\,\dd W_t^{\perp}.
  \label{eq:expression-sde}
\end{align}
The scalar clock rate and variability are state-dependent,
\begin{equation}
  a_\phi(y)=\operatorname{softplus}\!\left(h_a(E(Qy))\right),
  \qquad
  \sigma_\phi(y)=\operatorname{softplus}\!\left(h_\sigma(E(Qy))\right),
  \label{eq:clock-heads}
\end{equation}
where $h_a$ and $h_\sigma$ are small heads on frozen \uce{} state. The
direction $u_\theta$ retains the gene-specific CRE-conditioned architecture.
The same scalar realization $\dd S_t$ multiplies every gene for one cell, so
clock variation moves that cell coherently along its local developmental
direction. $G_\psi\dd W_t^\perp$ represents residual stochasticity such as
unresolved fate choice or technical variation.

The clock is a latent dynamical variable, not a conditioning input. In
particular, neither $S_t$ nor absolute developmental time is supplied to
$u_\theta$, $a_\phi$, or $\sigma_\phi$. Every field is a function of the
current expression-derived state.

\subsection{Positive progress increments}

A Gaussian Euler increment for Equation~\ref{eq:clock-sde} can occasionally be
negative. Small backward fluctuations may be biologically acceptable, but a
first implementation should test a monotone clock using positive lognormal
increments:
\begin{equation}
  \Delta S_k
  = a_\phi(Y_k)\Delta t
    \exp\!\left[
      \sigma_\phi(Y_k)\sqrt{\Delta t}\,\epsilon_k
      -\frac{1}{2}\sigma_\phi(Y_k)^2\Delta t
    \right],
  \qquad \epsilon_k\sim\N(0,1).
  \label{eq:positive-increment}
\end{equation}
The correction in the exponent preserves
$\E[\Delta S_k\mid Y_k]=a_\phi(Y_k)\Delta t$. The model-coordinate update is
\begin{equation}
  Y_{k+1}=\Pi_+\!\left[
    Y_k+u_\theta(Y_k)\Delta S_k
    +G_\psi(Y_k)\sqrt{\Delta t}\,\epsilon_k^\perp
  \right],
  \label{eq:clock-rollout}
\end{equation}
where $\Pi_+$ projects scaled expression to the nonnegative orthant. Count-like
expression is recovered exactly as $X_{k+1}=QY_{k+1}$ before \uce{} or
evaluation. This reparameterization is differentiable with respect to the
model parameters.

Gamma increments or an integrated positive stochastic rate can be considered
later. The lognormal increment is simple, has a direct
coefficient-of-variation interpretation, and requires one additional Gaussian
draw per cell and step.

\subsection{Direction normalization}

Without a constraint, multiplying $u_\theta$ by a constant and dividing
$a_\phi$ by the same constant leaves the expression dynamics unchanged. To
separate direction from progress, predict velocity directly in $Y$ coordinates
and use an ordinary root-mean-square normalization:
\begin{equation}
  u_\theta(y)
  = \frac{v_\theta(y)}
      {\sqrt{G^{-1}\sum_{g=1}^{G}v_{\theta,g}(y)^2}+\epsilon},
  \label{eq:direction-normalization}
\end{equation}
The clock increment then determines distance traveled in scaled expression
space. A less invasive
first experiment can freeze the pretrained $v_\theta$ and learn only the two
clock heads; Equation~\ref{eq:direction-normalization} is preferable when the
direction and clock are later trained jointly.

\subsection{Separating clock and residual noise}

Clock noise is rank one locally: its model-coordinate direction is
$u_\theta(y)$. Residual diffusion should not be allowed to relearn the same
component. Given an unconstrained residual draw $r$, project it away from the
scaled velocity direction,
\begin{equation}
  r_\perp
  = r-u\frac{r^Tu}{u^Tu}.
  \label{eq:orthogonal-noise}
\end{equation}
This decomposition makes the interpretation testable: $\sigma_\phi$ controls
variation in progress, while $G_\psi$ controls variation transverse to the
trajectory. The first ablation should omit $G_\psi$ entirely so that any gain
can be attributed to the clock.

\section{Learning without pseudotime labels}

\subsection{Unpaired population objective}

For each observed interval $(t_0,t_1)$, independently sample source cells
$\{x_i^0\}_{i=1}^{B}$ and target cells $\{x_j^1\}_{j=1}^{C}$. Roll $M$
particles from each source according to Equation~\ref{eq:clock-rollout}. No
cell-level OT coupling is constructed. Compare the generated endpoint
model-coordinate population $\widehat{\mathcal Y}_1$ with the observed target
population $\mathcal X_1$ after deterministic conversion back to expression:
\begin{align}
  \mathcal L_{\mathrm{endpoint}}
  =&\;S_\varepsilon\!\left(
      P(Q\widehat{\mathcal Y}_1),P(\mathcal X_1)
    \right)
  +\lambda_m\mathcal L_{\mathrm{mean}}
  +\lambda_C\mathcal L_{\mathrm{cov}},
  \label{eq:endpoint-loss}
\end{align}
where $P$ is a differentiable, train-only log-normalized PCA map and
$S_\varepsilon$ is debiased Sinkhorn divergence. Gene means are compared after
robust scale normalization, and covariance is compared in PCA space. This is
the same population-level supervision used by the explicit score-flow
fine-tuning implementation.

The experimental interval length enters only through the number and size of
integration steps. The model learns how much biological progress usually
occurs during that duration from endpoint distributions. It never receives a
precomputed pseudotime value for an individual cell.

\subsection{Clock-aligned flow and score matching}

The clock can also be pretrained without differentiating through full
rollouts. Standard flow matching by itself learns only the experimental-time
expression field
\begin{equation}
  b(x)=a_\phi(x)u_\theta(x).
  \label{eq:rate-direction-product}
\end{equation}
It cannot identify the rate and direction separately: multiplying
$u_\theta$ by a constant and dividing $a_\phi$ by the same constant leaves
$b$ unchanged. Direction normalization in
Equation~\ref{eq:direction-normalization}, a rate anchor, and supervision
across multiple interval lengths are therefore required before interpreting
$a_\phi$ as latent progress rate.

Given coupled expression endpoints, first set $y_i=Q^{-1}x_i$, then define the
model-coordinate chord length and direction
\begin{equation}
  d=y_1-y_0,\qquad
  \ell=\sqrt{G^{-1}\sum_gd_g^2},\qquad
  \widehat d=d/(\ell+\epsilon).
  \label{eq:clock-chord}
\end{equation}
A clock-aligned stochastic interpolant perturbs the linear path using one
scalar Gaussian variable shared by every gene:
\begin{align}
  Y_\tau
  &=(1-\tau)y_0+\tau y_1+\rho\gamma(\tau)\widehat d\,z,
  &z&\sim\N(0,1),
  \label{eq:clock-interpolant}\\
  \partial_\tau Y_\tau
  &=d+\rho\gamma'(\tau)\widehat d\,z.
  \label{eq:clock-flow-target}
\end{align}
Unlike the current gene-wise Gaussian interpolant, the random term changes
distance along the endpoint chord rather than independently perturbing every
gene. Mapping this path through $Q$ exactly gives a count-space linear
interpolant with gene-scaled noise. Here $\rho>0$ is a fixed dimensionless
bridge-noise amplitude, analogous
to the fixed noise amplitude in the current stochastic interpolant. The model
predicts a normalized direction $u_\theta(Y_\tau)$, positive rate
$a_\phi(Y_\tau)$, and scalar noise
$\widehat z_\psi(Y_\tau,\tau)\in\R$, all conditioned on $E(QY_\tau)$. The
rollout clock scale $\sigma_\phi(Y_\tau)$ is subsequently calibrated and fine-tuned by the
endpoint objective. A coupled clock-flow readout has the form
\begin{equation}
  f_{\mathrm{clock}}
  =a_\phi u_\theta
   +\frac{\gamma'(\tau)}{t_1-t_0}
    \rho u_\theta\widehat z_\psi,
  \label{eq:coupled-clock-flow}
\end{equation}
while the score component tangent to the trajectory is constructed from the
same scalar estimate. Since $G^{-1}u^Tu=1$, define the Euclidean dual direction
\begin{equation}
  u^*=\frac{u}{G},
  \qquad
  u^{*T}u=1.
  \label{eq:dual-direction}
\end{equation}
The tangent ambient-score vector is then
\begin{equation}
  s_\parallel
  =-\frac{\widehat z_\psi}
          {\rho\gamma(\tau)+\epsilon}\,u^*.
  \label{eq:tangent-clock-score}
\end{equation}
Thus the flow and score cannot disagree about the latent clock fluctuation.
Because rank-one clock noise is singular in the full gene space,
Equation~\ref{eq:tangent-clock-score} is a tangent score rather than a complete
ambient score. A small fixed residual variance or the orthogonal diffusion in
Equation~\ref{eq:orthogonal-noise} is needed if a nonsingular full-space score
is required.

An initial simulation-free objective can combine
\begin{equation}
  \mathcal L_{\mathrm{clock\mbox{-}FM}}
  =\mathcal L_{\mathrm{flow}}
   +\lambda_z\lVert\widehat z_\psi-z\rVert^2
   +\lambda_d\left[1-\cos(u_\theta,\widehat d)\right]
   +\lambda_a(\log a_\phi)^2.
  \label{eq:clock-fm-loss}
\end{equation}
The direction term should be downweighted or omitted once the model is
fine-tuned against population endpoints, because an individual endpoint chord
is a coupling-dependent training device rather than a known lineage.

Simulation-free stochastic-interpolant training still requires an endpoint
coupling. A large-pool partial OT or KL-relaxed UOT plan can supply diverse
training paths without interpreting pairs as true descendants. Avoiding a
fixed coupling entirely requires either the explicit unpaired rollout loss in
Equation~\ref{eq:endpoint-loss} or an iterative Schr\"odinger-bridge procedure
that alternates SDE simulation and score fitting. The latter remains a
flow/score method, but is more expensive and introduces an iterative bridge
estimation loop.

Gaussian clock score matching corresponds naturally to
$\dd S=a\,\dd t+\sigma\,\dd W$ and therefore permits occasional backward
increments. The strictly positive lognormal increment in
Equation~\ref{eq:positive-increment} does not have the same simple Gaussian
denoising target. Consequently, the recommended hybrid is:
\begin{enumerate}
  \item pretrain normalized direction, rate, and scalar clock score with
        clock-aligned stochastic-interpolant matching;
  \item freeze or strongly anchor the normalized direction; and
  \item fine-tune the positive stochastic clock with the unpaired endpoint
        rollout objective.
\end{enumerate}
This uses flow/score matching for efficient local field estimation and the
population objective for monotonicity and coupling-independent model
selection.

\subsection{Anchors and regularization}

The endpoint objective alone admits weakly identified solutions. The initial
loss should therefore include
\begin{equation}
  \mathcal L
  =\mathcal L_{\mathrm{endpoint}}
   +\lambda_a\E[(\log a_\phi)^2]
   +\lambda_\sigma\E[\sigma_\phi^2]
   +\lambda_\theta\lVert\theta-\theta_0\rVert_2^2.
  \label{eq:regularized-loss}
\end{equation}
The rate prior centers $a_\phi$ near one in the chosen progress units, the
noise penalty prevents variance inflation from substituting for dynamics, and
$\theta_0$ is the simulation-free pretrained velocity checkpoint. These are
soft anchors rather than pseudotime labels.

Validation must use cells disjoint from training within every observed stage,
fixed source and target populations, fixed clock-noise draws, and an
interval-balanced average. The pretrained fixed-clock model is validation
step zero. Checkpoint selection should require a material improvement over
that baseline rather than merely a decreasing training objective.

\subsection{Truncated gradients through UCE}

At each rollout step, \uce{} is recomputed from current raw expression using
the low-variance systematic tokenization scheme. Tokenization is discrete and
the encoder is frozen, so its state should be treated as stop-gradient.
Gradients still reach every clock and velocity evaluation directly and pass
through the additive expression updates. KNN or metacell smoothing may be used
to construct simulation-free pretraining targets, but it should not be applied
during causal clock rollout because a same-stage reference population is not
available for a future cell.

\section{Interpretation and limitations}

\subsection{What the clock can explain}

The model is appropriate when cells share a local developmental path but vary
in how quickly they move along it. It yields interpretable quantities:
\begin{itemize}
  \item $a_\phi(x)$: expected progress rate for a cell state;
  \item $\sigma_\phi(x)$: uncertainty or heterogeneity in that rate;
  \item $S_{t_1}-S_{t_0}$: realized latent progress over an interval; and
  \item $u_\theta(x)$: gene-expression direction per unit progress.
\end{itemize}
Stage-wise distributions of inferred progress should be ordered on average,
but overlap is expected and is one motivation for the model.

\subsection{Branching}

A scalar clock does not resolve branching when the same expression state can
enter qualitatively different futures. Clock noise changes distance along a
direction, not the direction itself. If the clock improves short-range
forecasting but fails near fate decisions, the next extension should use a
small mixture of regulatory velocity directions with a state-dependent gate,
or retain an explicitly low-rank transverse diffusion. Branch-specific
complexity should not be added until the rank-one clock ablation establishes
that variable progress is useful.

\subsection{Proliferation and changing population mass}

Neither the clock nor a mass-preserving Sinkhorn loss explains changes in cell
type abundance caused by division, death, or sampling. A progression-rate
head may correlate with proliferation, but it does not create or remove
particles. Growth weights, a learned birth--death rate, or unbalanced endpoint
loss remain complementary. This distinction is important: faster progress,
trajectory branching, and population growth are separate mechanisms and
should not be absorbed into one noise parameter.

\section{Ablation plan}

\begin{table}[ht]
\centering
\small
\begin{tabular}{@{}p{0.06\textwidth}p{0.25\textwidth}p{0.24\textwidth}p{0.27\textwidth}@{}}
\toprule
Arm & Dynamics & Learned additions & Question \\
\midrule
A & $\dd Y=v_\theta\dd t$ & None & Fixed-clock persistence-relative baseline \\
B & $\dd Y=a_\phi u_\theta\dd t$ & Deterministic rate & Does state-dependent speed help? \\
C & $\dd Y=u_\theta\dd S$ & Rate and scalar clock noise & Does variable progress improve endpoint prediction? \\
D & Current score SDE & Gene-wise score/noise & Is generic diffusion sufficient? \\
E & Clock plus $G_\psi\dd W^\perp$ & Clock and residual noise & Is transverse stochasticity additionally required? \\
\bottomrule
\end{tabular}
\caption{Staged ablations separating variable developmental progress from generic expression diffusion.}
\label{tab:ablations}
\end{table}

Arms A--C should be implemented first, using identical independent endpoint
batches and Brownian seeds. Begin by freezing the pretrained regulatory
velocity, then repeat B and C with normalized direction fine-tuning if either
shows held-out benefit. Arm E is justified only if C improves progress-related
metrics but leaves systematic multimodal endpoint errors.

\section{Evaluation}

The primary endpoint metric remains persistence-normalized Sinkhorn skill,
\begin{equation}
  \operatorname{skill}
  =1-\frac{S_\varepsilon(P(\widehat{\mathcal X}_1),P(\mathcal X_1))}
           {S_\varepsilon(P(\mathcal X_0),P(\mathcal X_1))}.
  \label{eq:skill}
\end{equation}
Positive skill means the predicted population is closer to the observed
endpoint than a no-change forecast. This must be accompanied by gene-level
mean-shift $R^2$, mean-shift MAE, per-gene Wasserstein distance, and major-cell-
type-stratified metrics. Sinkhorn improvement alone can reward variance
broadening.

Clock-specific diagnostics should include:
\begin{itemize}
  \item distributions of $a_\phi$, $\sigma_\phi$, and realized $\Delta S$ by
        experimental stage and major cell type;
  \item the fraction of endpoint variance explained by clock-aligned versus
        transverse displacement;
  \item stability of rates under Brownian seed, particle count, and rollout
        step count;
  \item calibration of predicted progress variance against observed spread
        along the local velocity direction; and
  \item correlations with cell-cycle scores as diagnostics, not training
        labels, to detect a clock that merely recovers proliferation state.
\end{itemize}

Success requires a held-out-stage gain over the fixed-clock model, not only an
observed-interval validation gain. The learned clock should improve both
distributional and gene-shift metrics, remain stable across integration
settings, and produce bounded state-dependent rates rather than extreme
variance in a small number of cells.

\section{Implementation plan}

The implementation should preserve the existing separation between reusable
model code and atlas-specific scripts:
\begin{enumerate}
  \item Fit $q$ from raw training cells only and add a checkpointed
        $Y=Q^{-1}X$ coordinate transform. Keep OT/PCA and \uce{} in $X$, but
        perform interpolation, velocity loss, score matching, and rollout in
        $Y$. Preserve legacy count-space checkpoint contracts.
  \item Add reusable scalar rate and clock-noise heads operating on frozen
        \uce{} state. Initialize the deterministic rate near one and clock
        noise near zero.
  \item Add the positive stochastic increment in
        Equation~\ref{eq:positive-increment} to the differentiable population
        rollout module. Accept fixed scalar Gaussian draws for reproducible
        validation.
  \item Implement robust direction normalization and verify invariance to a
        global rescaling of raw velocity.
  \item Add the scalar clock-aligned stochastic interpolant and coupled
        flow/noise loss in Equations~\ref{eq:clock-interpolant}
        and~\ref{eq:clock-fm-loss}. Pretrain with matched large-pool partial-OT
        or UOT paths.
  \item Fit Arm B and Arm C from the same pretrained score-flow or autonomous
        checkpoint, initially freezing the regulatory direction, then
        fine-tune with unpaired positive-clock rollouts.
  \item Extend checkpoints with clock configuration, initialization,
        normalization scales, parent checkpoint, and selected validation step.
  \item Extend held-out causal evaluation with fixed-clock, deterministic-rate,
        and stochastic-clock controls using identical source cells.
  \item Add synthetic tests for known variable-speed trajectories, including a
        case where diagonal diffusion broadens the population but cannot match
        coherent progress variation.
  \item Run matched raw, KNN-pretrained, and metacell-pretrained models without
        applying smoothing during rollout.
\end{enumerate}

The first experiment should be small: two particles per source, two to four
rollout steps, a deterministic-rate arm, and a stochastic-clock arm. The
existing explicit population fine-tuning workflow already supplies independent
endpoint batches, differentiable PCA/Sinkhorn loss, frozen validation draws,
and pretrained checkpoint anchoring. The main new machinery is therefore two
scalar heads and one shared random increment per cell and step.

\section{Decision rule}

The stochastic clock is supported if Arm C improves held-out endpoint skill
and gene-shift metrics over both Arm A and Arm B, and if learned progress
variance is nonzero but bounded. If Arm B improves while Arm C does not, cells
benefit from state-dependent deterministic rates but stochastic pseudotime is
unnecessary. If the current gene-wise SDE improves while the clock does not,
the useful uncertainty is not primarily variation in developmental progress.
If no arm exceeds persistence, the limiting issue is more likely trajectory
direction, lineage ambiguity, growth, or insufficient information in snapshot
expression than the absence of a pseudotime coordinate.

\begin{thebibliography}{99}

\bibitem{haghverdi2016dpt}
Haghverdi, L., B\"uttner, M., Wolf, F. A., Buettner, F. \& Theis, F. J.
Diffusion pseudotime robustly reconstructs lineage branching.
\textit{Nature Methods} \textbf{13}, 845--848 (2016).

\bibitem{setty2019palantir}
Setty, M. et al.
Characterization of cell fate probabilities in single-cell data with Palantir.
\textit{Nature Biotechnology} \textbf{37}, 451--460 (2019).

\bibitem{bergen2020scvelo}
Bergen, V., Lange, M., Peidli, S., Wolf, F. A. \& Theis, F. J.
Generalizing RNA velocity to transient cell states through dynamical modeling.
\textit{Nature Biotechnology} \textbf{38}, 1408--1414 (2020).

\bibitem{schiebinger2019wot}
Schiebinger, G. et al.
Optimal-transport analysis of single-cell gene expression identifies
developmental trajectories in reprogramming.
\textit{Cell} \textbf{176}, 928--943 (2019).

\bibitem{song2021sde}
Song, Y. et al.
Score-Based Generative Modeling through Stochastic Differential Equations.
\textit{International Conference on Learning Representations} (2021).
\url{https://arxiv.org/abs/2011.13456}

\bibitem{li2020neuralsde}
Li, X., Wong, T.-K. L., Chen, R. T. Q. \& Duvenaud, D.
Scalable Gradients for Stochastic Differential Equations.
\textit{International Conference on Artificial Intelligence and Statistics}
(2020).
\url{https://arxiv.org/abs/2001.01328}

\bibitem{albergo2023stochastic}
Albergo, M. S., Boffi, N. M. \& Vanden-Eijnden, E.
Stochastic Interpolants: A Unifying Framework for Flows and Diffusions.
\url{https://arxiv.org/abs/2303.08797}

\end{thebibliography}

\end{document}
